\chapter{\abstractname}
Medical imaging repositories are large, heterogeneous, and operationally complex. In research and quality assurance (QA) workflows, scalable access to Digital Imaging and Communications in Medicine (DICOM) metadata is often more important than pixel rendering, yet typical tooling remains fragmented between ad hoc scripts, file-level utilities, and heavyweight enterprise systems. This thesis presents a complete DICOM Metadata Browser that transforms heterogeneous DICOM file collections into compact, queryable databanks with integrated quality analytics.

The implemented system combines robust DICOM discovery, metadata-only extraction, preservation of private tags, selection of representative series, and an interactive web user interface (UI). It supports mixed extension conventions (\texttt{.dcm}, \texttt{.ima}, extensionless files), nested scan directories, single-file ingestion, ZIP archive inputs, deduplication, and multilingual browsing (English and German). The storage layer uses SQLite with indexing and performance-oriented settings, augmented by a dedicated private-tag table that preserves creator information, value fingerprints, and classifications. Siemens Common Syngo Architecture (CSA) payloads are parsed and persisted as structured JSON with hashes for consistency checks.

Evaluation on four local datasets indicates fast local processing and strong storage reduction. Measured end-to-end processing times range from 1.44\,s to 13.14\,s in the test environment. Database size reduction compared with the processed scan folders ranges from 19.2x to 147.3x. Integrated quality checks identify completeness gaps and temporal anomalies, including negative injection-to-acquisition delays in one subset. Taken together, these results suggest that a local-first, pragmatically engineered metadata platform can provide meaningful operational value for imaging research and QA without requiring heavy infrastructure. 
